\paragraph{Risposta (b)}
\begin{itemize}
\item
  Se per ipotesi il filtro \'e in regione di linearit\'a allora $V_+ = V_-$ e siccome
  $V_+ = 0$ \'e messo a terra lo sar\'a anche $V_- = 0$.
\end{itemize}
La tensione all'ingresso dell'invertente \'e il risultato di una partizione di tensione
sul nodo A, quindi:
\begin{equation*}
  V_- = \frac{R_3}{R_3 + \frac{1}{sC_1}}V_A + \frac{\frac{1}{sC_1}}{R_3 + \frac{1}{sC_1}}V_{out}
\end{equation*}
per le considerazioni fatte in precedenza sull'idealit\'a dell'opamp
\begin{equation*}
  0 = \frac{R_3}{R_3 + \frac{1}{sC_1}}V_A + \frac{\frac{1}{sC_1}}{R_3 + \frac{1}{sC_1}}V_{out} \\
  V_{out} = -R_3C_1V_A
\end{equation*}
Applicando la legge di Kirchoff sul nodo A
\begin{equation*}
  \frac{V_s - V_A}{R_1} = \frac{V_A - V_{out}}{\frac{1}{sC_2}} + \frac{V_A - V_-}{\frac{1}{sC_1}} + \frac{V_A - V_+}{R_2}
\end{equation*}
per le considerazioni fatte in precedenza sull'idealit\'a dell'opamp
\begin{equation*}
  \frac{V_s - V_A}{R_1} = \frac{V_A - V_{out}}{\frac{1}{sC_1}} + \frac{V_A}{\frac{1}{sC_1}} + \frac{V_A}{R_2}
\end{equation*}
